%%%%%%%%%%%%%%%%%%%%%%%%%%%%%%%%%%%%%%%%%%%%%%%%%%%%%%%%%%%%%%%%%%%%%%%%%%%
%%  Project: CPHOS LaTeX Template
%%  File: example-paper.tex
%%  Copyright © 2026 gameswu & CPHOS (www.cphos.cn)
%%
%%  Permission is hereby granted, free of charge, to any person obtaining a 
%%  copy of this software and associated documentation files (the "Software"), 
%%  to deal in the Software without restriction, including without limitation 
%%  the rights to use, copy, modify, merge, publish, distribute, sublicense, 
%%  and/or sell copies of the Software, and to permit persons to whom the 
%%  Software is furnished to do so, subject to the following conditions:
%%
%%  The above copyright notice and this permission notice shall be included 
%%  in all copies or substantial portions of the Software.
%%
%%  THE SOFTWARE IS PROVIDED "AS IS", WITHOUT WARRANTY OF ANY KIND, EXPRESS 
%%  OR IMPLIED, INCLUDING BUT NOT LIMITED TO THE WARRANTIES OF MERCHANTABILITY, 
%%  FITNESS FOR A PARTICULAR PURPOSE AND NONINFRINGEMENT. IN NO EVENT SHALL 
%%  THE AUTHORS OR COPYRIGHT HOLDERS BE LIABLE FOR ANY CLAIM, DAMAGES OR OTHER 
%%  LIABILITY, WHETHER IN AN ACTION OF CONTRACT, TORT OR OTHERWISE, ARISING 
%%  FROM, OUT OF OR IN CONNECTION WITH THE SOFTWARE OR THE USE OR OTHER 
%%  DEALINGS IN THE SOFTWARE.
%%%%%%%%%%%%%%%%%%%%%%%%%%%%%%%%%%%%%%%%%%%%%%%%%%%%%%%%%%%%%%%%%%%%%%%%%%%


\documentclass[exam]{cphos-e}

\cphostitle{第XX届 CPHOS 物理竞赛联考}
\cphossubtitle{实验试题}
\cphoswatermark{../assets/watermark.jpg}

\begin{document}

% === 标题页 ===
\maketitlepage

{\kaiti\small 本文件于YYYY年MM月DD日HH:MM发布，最后更新于\compiletime 。}

{\kaiti\small CPHOS物理竞赛联考是开放性公益性的考试，有意向参与的教师和学生可以关注“CPHOS”微信公众号进行报名，报名后方可参与联考。请使用“CPHOS物理竞赛联考”微信小程序完成答题卡上传、阅卷、成绩查询等操作。联系方式见试题末尾。}
% === 时间安排（仅exam模式显示）===
\begin{examschedule}
\begin{tabular}{@{}cc@{}}
    答题卡上传 & 阅卷 \\
    YYYY/MM/DD HH:MM - YYYY/MM/DD HH:MM & YYYY/MM/DD HH:MM - YYYY/MM/DD HH:MM \\
\end{tabular}

\vspace{0.3em}

\begin{tabular}{@{}ccc@{}}
    非正式成绩 & 成绩申诉 & 正式成绩 \\
    YYYY/MM/DD HH:MM & YYYY/MM/DD HH:MM - YYYY/MM/DD HH:MM & YYYY/MM/DD HH:MM \\
\end{tabular}
\end{examschedule}

% === 考生须知（仅exam模式显示）===
\begin{examnotes}
    \item 实验试题共 \textbf{\underline{\cphostotalpages}} 页，实验答题卡共 \textbf{\underline{X}} 页，答题时间 \textbf{\underline{60}} 分钟，试题满分 \textbf{\underline{\cphostotalscore}} 分。
    \item 请在答题卡的指定答题区域内答题，试题和草稿纸上的内容将不会作为评分参考，不可申请答题卡加页。
    \item 若发现试题存在问题，请向领队（教练）反映，由其转达至相关微信群聊。
    \item 试题答案及相关分析均会在官方网站 \url{www.cphos.cn} 上发布。
    \item 本次考试定位难度为决赛。
\end{examnotes}
\vspace{1em}

\begin{problem}[40]{弗兰克-赫兹实验}
\begin{expcontent}

1913 年玻尔模型指出原子能级是分立的。弗兰克-赫兹实验通过电子与气体原子碰撞电离/激发行为，验证了原子能级量子化。

\begin{figure}[H]
    \centering
    \fbox{\parbox[c][5.5cm][c]{0.72\textwidth}{\centering 实验装置图占位示意}}
    \caption{弗兰克-赫兹实验装置示意图}
    \label{fig:fh-device}
\end{figure}

\pmark{观察与定性分析}
\subq 在固定灯丝加热电压与栅极电压条件下，调节加速电压并记录电流变化特征，说明首次电流峰值出现前后电子与汞原子碰撞行为的变化。
\subq 讨论异常下降段可能对应的实验误差来源，并给出修正建议。

\pmark[15]{数据处理与参数估计}
\subq 根据实验表格中相邻峰值间隔，估算第一激发能对应电势差。
\subq 结合误差传播，给出测量结果的不确定度表达。
\subq 比较实验值与参考值，说明是否在误差范围内一致。
\end{expcontent}
\end{problem}

% ============================================
% 版权信息（使用copyrightinfo环境确保不分页）
% ============================================
\begin{copyrightinfo}
\begin{center}
    {\textbf{版权信息}}
\end{center}

\noindent{\textbf{命题人}}

X~X\quad XXX\quad XXX

\noindent{\textbf{审题人}}

Y~Y\quad YYY\quad YYY

\vspace{0.5em}
\contactinfo{../assets/fig1.jpg}{../assets/fig2.png}{../assets/fig3.png}{\url{service@cphos.cn}}{CPHOS物理竞赛联考}
\end{copyrightinfo}

\end{document}
